\subsection{在线归一化同步核的 \texorpdfstring{$\zeta_B$}{zeta\_B} 与 primitive 轨道谱}\label{subsec:zeta-online-kernel}
第 \ref{subsec:online-delay} 小节的单趟在线延迟转导器，其同步部分给出一个 10 状态 SFT（见附录 B），记邻接矩阵为 $B$。按 SFT 的标准公式，
\[
\zeta_B(z)=\frac{1}{\det(I-zB)},\qquad
a_n:=\#\mathrm{Fix}(\sigma^n)=\mathrm{Tr}(B^n).
\]
固定状态顺序与附录 B 一致，并令
\[
B_{ij}:=\#\{d\in\{0,1,2\}:\ q_i\xrightarrow{d}q_j\},
\]
则由同步转移表逐边计数即可得到该 $B$；其显式矩阵与计算细节见附录 B。
对该同步核可得显式因式分解（附录 B 给出矩阵与计算）
\[
\det(I-zB)=(z-1)(z+1)(3z-1)\bigl(z^3-z^2+z+1\bigr),
\]
因此主极点位于 $z_\star=1/3$，对应 Perron--Frobenius 主特征值 $\lambda=3$，并且
\[
\zeta_B(z)=\frac{C_B}{1-3z}+O(1),\qquad C_B=\frac{243}{272}
\]
（$C_B$ 为该极点留数，见附录 B）。令 $p_n$ 为长度 $n$ 的 primitive 轨道数（按轨道计），则
\[
a_n=\sum_{d\mid n} d\,p_d,
\qquad
p_n=\frac{1}{n}\sum_{d\mid n}\mu(d)\,a_{n/d}.
\]
由 $B$ 的原始性（例如 $B^5$ 全正），得到
\[
a_n=3^n+O(\rho^n),\qquad
p_n=\frac{3^n}{n}+O\!\left(\frac{\rho^n}{n}\right)\quad(\rho<3).
\]
对应的“素数计数曲线”定义为 $\Pi_B(N):=\sum_{n\le N}p_n$，其离散主项为
\[
\Pi_B(N)\sim \frac{3}{2}\cdot \frac{3^N}{N},
\]
并可写成 $X=3^N$ 的 $X/\log X$ 形态（详见附录 B）。可复现实验脚本见 \texttt{scripts/exp\_sync\_kernel\_primitive\_spectrum.py}。

\paragraph{加权 $\zeta(z,u)$ 的函数方程与输出密度常数}
取输出位势 $\varphi(e)=e$，令 $B(u)=B_0+uB_1$、$\Delta(z,u)=\det(I-zB(u))$ 与 $\zeta(z,u)=\Delta(z,u)^{-1}$（见附录 \ref{app:witt-weighted}）。附录 \ref{app:kernel-self-dual} 给出自对偶对称
\[
\Delta(z,u)=\Delta(uz,1/u),\qquad
\zeta(z,u)=\zeta(uz,1/u),
\]
从而 $P(\theta)=\log\lambda(e^\theta)$ 满足 $P(\theta)=\theta+P(-\theta)$。同理令
\[
\widetilde{\Delta}(z,\theta):=\Delta(e^{-\theta/2}z,e^\theta),
\]
则完成化行列式在 $\theta$ 上严格偶对称（见附录 \ref{app:kernel-completion}）。因此
\[
P'(0)=\frac12,\qquad P''(0)=\frac{11}{102},
\]
并且除均值外的所有奇阶累积量密度皆为零（$P^{(2k+1)}(0)=0,\ k\ge 1$，见附录 \ref{app:pressure-analytic}）。
并在典型密度 $\alpha_\star=\frac12$ 附近有
\[
I(\alpha)=\frac{(\alpha-\frac12)^2}{2P''(0)}+O\!\left((\alpha-\tfrac12)^3\right)
=\frac{51}{11}\Bigl(\alpha-\frac12\Bigr)^2+O(\cdot),
\]
其中 $I$ 为输出位计数的 LDP 率函数（见附录 \ref{app:pressure-analytic} 与 \ref{app:rate-param}）。有限 $n$ 的镜像对称与其导出的 Mertens 偶性见推论 \ref{cor:primitive-finite-sym} 与命题 \ref{prop:mertens-even}；周期迹多项式与 $\Delta$ 系数的逐项回文约束见命题 \ref{prop:trace-palindrome} 与推论 \ref{cor:delta-coeff-palindrome}。

\paragraph{Cyclotomic elimination：模 \texorpdfstring{$m$}{m} 的 residue mixing 误差指数是整系数代数数}
把加权 primitive 轨道谱写成
\[
p_n(u)=\sum_{k}p_{n,k}\,u^k\in\ZZ[u],
\]
其中 $k$ 是长度 $n$ 的 primitive 轨道上输出 $1$ 的出现次数。对任意模 $m\ge 3$ 与同余类 $a\in\ZZ/m\ZZ$，定义 residue class 的 primitive 计数
\[
p_{n,a}^{(m)}:=\sum_{k\equiv a\ (m)}p_{n,k}.
\]
由离散傅里叶反演（evaluation@roots + DFT）有严格恒等式
\[
p_{n,a}^{(m)}=\frac{1}{m}\sum_{r=0}^{m-1}\omega_m^{-ar}\,p_n(\omega_m^r),
\qquad \omega_m:=e^{2\pi i/m}.
\]
因此 residue mixing 的误差指数由扭曲 Perron 根
\[
\rho_m:=\max_{1\le r\le m-1}|\lambda(\omega_m^r)|<\lambda(1)=3
\]
控制（$r=0$ 为平凡频率给出主项）。下面把 $\rho_m$ 变成一个\emph{整系数代数方程}（resultant）问题。

\begin{proposition}[completion \texorpdfstring{$\to$}{->} cyclotomic resultant]\label{prop:sync-cyclotomic-elimination}
记 $\Delta(z,u)=\det(I-zB(u))$。引入完成化变量
\[
u=r^2,\qquad s:=r+r^{-1},\qquad w:=zr,
\]
则由自对偶对称 $\Delta(z,u)=\Delta(uz,1/u)$ 可知 $w$ 与 $s$ 为不变量。于是存在
\[
\widehat{\Delta}(w,s)\in\ZZ[s][w]
\]
使得对所有可逆 $r$ 都有恒等式
\[
\widehat{\Delta}(w,r+r^{-1})=\Delta(w/r,r^2).
\]
令 $\Phi_m^+(s)$ 为代数数 $s_m:=2\cos(\pi/m)$ 的最小多项式（整系数、首一），并定义
\[
R_m(w):=\operatorname{Res}_s\bigl(\widehat{\Delta}(w,s),\Phi_m^+(s)\bigr)\in\ZZ[w].
\]
则当 $u=e^{2\pi i/m}$ 时，$\Delta(z,u)=0$ 的解 $z$（等价地 $w=z\sqrt{u}$）必满足 $R_m(w)=0$。并且在单位圆扭曲 $|u|=1$ 下有
\[
\rho_m=\frac{1}{\min\{|w|:R_m(w)=0\}}.
\]
\end{proposition}

\begin{remark}[可复现输出：\texorpdfstring{$m=3,\dots,7$}{m=3..7} 的 \texorpdfstring{$R_m$}{Rm} 与 \texorpdfstring{$\rho_m$}{rhom}]\label{rem:sync-cyclotomic-elimination-table}
脚本 \texttt{scripts/exp\_sync\_kernel\_cyclotomic\_elimination.py} 由附录 \ref{app:pressure-analytic} 的 $F(\lambda,u)$ 构造 $\widehat{\Delta}(w,s)$，并对 $m=3,4,5,6,7$ 计算 $R_m(w)$ 与 $\rho_m$；输出见表 \ref{tab:sync_kernel_cyclotomic_elimination_summary} 与多项式块 \texttt{tab\_sync\_kernel\_cyclotomic\_elimination\_polys.tex}（可直接回归测试）。
\end{remark}

\begin{table}[H]
\centering
\scriptsize
\setlength{\tabcolsep}{6pt}
\caption{Cyclotomic elimination for mod-$m$ prime-shadow mixing exponents of the weighted sync-kernel. $R_m(w)$ is an integer resultant polynomial; $\rho_m$ is obtained as $\rho_m=1/\min\{|w|:R_m(w)=0\}$. We also list the pressure-based prediction $\rho_m^{\mathrm{pred}}$ from the Taylor expansion of $P(\theta)$.}
\label{tab:sync_kernel_cyclotomic_elimination_summary}
\begin{tabular}{r r r r r}
\toprule
$m$ & $\deg R_m$ & $\rho_m$ (resultant) & $\rho_m^{\mathrm{pred}}$ & abs err\\
\midrule
3 & 5 & 2.418034345100 & 2.418835165986 & 8.008e-04\\
4 & 12 & 2.644080451866 & 2.643950096830 & 1.304e-04\\
5 & 12 & 2.762782720637 & 2.762705118214 & 7.760e-05\\
6 & 12 & 2.831532258235 & 2.831497400579 & 3.486e-05\\
7 & 18 & 2.874530783058 & 2.874514708145 & 1.607e-05\\
8 & 24 & 2.903081377378 & 2.903073483799 & 7.894e-06\\
9 & 18 & 2.922953899166 & 2.922949770585 & 4.129e-06\\
10 & 24 & 2.937319429068 & 2.937317144301 & 2.285e-06\\
11 & 30 & 2.948029997898 & 2.948028669926 & 1.328e-06\\
12 & 24 & 2.956223089850 & 2.956222284562 & 8.053e-07\\
\bottomrule
\end{tabular}
\end{table}

% AUTO-GENERATED by scripts/exp_sync_kernel_cyclotomic_elimination.py
\begin{align*}
R_{3}(w)&=5 w^{5} - 4 w^{4} - 3 w^{3} + 5 w^{2} + w - 1\\
R_{4}(w)&=w^{12} - 26 w^{10} + 47 w^{8} - 62 w^{6} + 43 w^{4} - 12 w^{2} + 1\\
R_{5}(w)&=w^{12} - 9 w^{11} + 13 w^{10} + 23 w^{9} - 33 w^{8} - 30 w^{7} + 52 w^{6} + 23 w^{5} - 38 w^{4} - 8 w^{3} + 11 w^{2} + w - 1\\
R_{6}(w)&=4 w^{12} - 19 w^{10} + 38 w^{8} - 61 w^{6} + 47 w^{4} - 13 w^{2} + 1\\
R_{7}(w)&=w^{18} + 11 w^{17} + 4 w^{16} - 107 w^{15} + 83 w^{14} + 249 w^{13} - 288 w^{12} - 352 w^{11} + 492 w^{10} + 338 w^{9} - 526 w^{8} - 208 w^{7} + 325 w^{6} + 75 w^{5} - 107 w^{4} - 14 w^{3} + 17 w^{2} + w - 1\\
\end{align*}


\begin{remark}[压力展开给出 \texorpdfstring{$\rho_m$}{rhom} 的高精度渐近预测]\label{rem:sync-rho-pressure-asymp}
由附录 \ref{app:pressure-analytic} 的 Taylor 展开
\[
P(\theta)=\log 3+\frac12\theta+\frac{11}{204}\theta^2+\frac{1559}{1414944}\theta^4+O(\theta^6),
\]
取 $\theta=it$ 并记 $u=e^{it}$，则
\[
\log|\lambda(e^{it})|=\Re P(it)=\log 3-\frac{11}{204}t^2+\frac{1559}{1414944}t^4+O(t^6).
\]
令 $t=2\pi/m$，得到可直接使用的预测式
\[
\rho_m\approx 3\exp\!\Bigl(-\frac{11}{204}\Bigl(\frac{2\pi}{m}\Bigr)^2+\frac{1559}{1414944}\Bigl(\frac{2\pi}{m}\Bigr)^4\Bigr),
\]
其与上面的 resultant 精确值在 $m=4$ 起已达到 $10^{-4}$ 量级误差（见表 \ref{tab:sync_kernel_cyclotomic_elimination_summary}）。
\end{remark}

\paragraph{\texorpdfstring{$s$}{s}-平面零极点字典：谱半径与“右半临界带无极点”}\label{par:s-plane-pole-dictionary}
本节的 \(\zeta\)-函数在有限矩阵口径下是有理函数，因此其全部奇点由有限谱完全决定。
为把“谱隙版 RH/GRH”提升到经典的 \((s)\)-平面语言，引入 Dirichlet 坐标
\[
z=\lambda^{-s},\qquad \lambda:=\rho(M)>0.
\]
\begin{proposition}[谱 \(\Longleftrightarrow\) Dirichlet 极点列]\label{prop:s-plane-spectrum-poles}
设 \(M\) 为有限维矩阵，\(\lambda:=\rho(M)>0\)。定义
\[
Z_M(s):=\frac{1}{\det\!\bigl(I-\lambda^{-s}M\bigr)}.
\]
则：
\begin{enumerate}
  \item \(Z_M(s)\) 的极点与 \(M\) 的非零特征值 \(\mu\) 一一对应（按代数重数）。更具体，\(\lambda^{-s}\mu=1\) 当且仅当
  \[
  s=\frac{\log\mu+2\pi i k}{\log\lambda}\qquad(k\in\ZZ),
  \]
  其中 \(\log\mu\) 为复对数的任一分支。
  \item 极点所在竖线的实部由模长给出：
  \[
  \Re(s)=\frac{\log|\mu|}{\log\lambda}.
  \]
\end{enumerate}
因此，核版 RH 条件
\(
\max_{i\ge 2}|\mu_i|\le \sqrt{\lambda}
\)
等价于：除主极点列（来自 \(\mu=\lambda\)，对应 \(\Re(s)=1\)）外，所有极点皆满足 \(\Re(s)\le \tfrac12\)，亦即 \(\Re(s)>\tfrac12\) 的区域无非平凡极点贡献。
\end{proposition}
\begin{proof}
对 \(\det(I-\lambda^{-s}M)\) 取 Jordan 分解或特征多项式分解即可：其零点满足 \(\lambda^{-s}\mu=1\)；取模得到 \(\lambda^{-\Re(s)}|\mu|=1\)，即 \(\Re(s)=\log|\mu|/\log\lambda\)。证毕。
\end{proof}

\paragraph{有限维显式公式：素数影子的误差即“临界圆谱贡献”}\label{par:finite-explicit-formula}
对有限维核，primitive 轨道计数的“主项 + 误差”可完全写成特征值的有限和，从而把平方根振荡直接识别为临界圆 \(|\mu|=\sqrt{\lambda}\) 上的谱贡献。
\begin{proposition}[有限维显式公式模板]\label{prop:finite-dimensional-explicit-formula}
设 \(M\) 的特征值（按代数重数）为 \(\mu_1,\dots,\mu_d\)，并记
\(
a_n:=\Tr(M^n)=\sum_{j=1}^d \mu_j^n
\)。
令 \(p_n\) 为长度 \(n\) 的 primitive 轨道数（按轨道计），则 Möbius 反演给出
\[
p_n=\frac{1}{n}\sum_{m\mid n}\mu(m)\,a_{n/m}
=\frac{1}{n}\sum_{j=1}^d\ \sum_{m\mid n}\mu(m)\,\mu_j^{\,n/m}.
\]
特别地，若 \(\lambda:=\rho(M)\) 为简单 Perron 根，则
\[
p_n=\frac{\lambda^n}{n}+O\!\left(\frac{\Lambda^n}{n}\right),
\qquad
\Lambda:=\max_{j\ge 2}|\mu_j|.
\]
若存在某个非平凡特征值满足 \(|\mu|=\sqrt{\lambda}\)，则一般出现不可改进的平方根级振荡项（\(\Omega(\lambda^{n/2}/n)\)）。
\end{proposition}
\begin{proof}
由 \(a_n=\sum_{m\mid n} m\,p_m\) 的标准周期--primitive 关系做 Möbius 反演得到第一式。
后两式由谱分解 \(a_n=\lambda^n+\sum_{j\ge 2}\mu_j^n\) 代回并估计即可；当某个 \(|\mu|=\sqrt{\lambda}\) 时，对应项规模为 \(\lambda^{n/2}\) 且无法被更小模的项一致抵消，从而给出 \(\Omega\)-振荡。证毕。
\end{proof}

\paragraph{\texorpdfstring{$s$}{s}-平面完成化函数方程的谱互反判据}\label{par:s-plane-functional-eq-criterion}
上一段的“函数方程”发生在参数侧（$u\leftrightarrow 1/u$ 或 $\theta\leftrightarrow-\theta$）。
若希望把对称进一步提升为 Dirichlet 坐标 $z=\lambda^{-s}$ 下的
\(
s\leftrightarrow 1-s
\)
型对称，则需要一个独立于参数自对偶的谱互反结构。

\begin{proposition}[谱互反闭包 $\Rightarrow$ 完成化后 \texorpdfstring{$s\leftrightarrow 1-s$}{s<->1-s}]\label{prop:reciprocal-spectrum-functional-equation}
设 $M$ 为有限维矩阵，$\lambda:=\rho(M)>0$ 为其 Perron 根。令
\[
\Delta_M(z):=\det(I-zM),\qquad
Z_M(s):=\Delta_M(\lambda^{-s})^{-1}.
\]
记 $r:=\deg_z\Delta_M(z)$（等价于 $M$ 的非零特征值按代数重数计的个数）。
若 $M$ 的非零谱满足\textbf{互反闭包}
\[
\mu\in\mathrm{spec}(M)\setminus\{0\}
\ \Longrightarrow\
\frac{\lambda}{\mu}\in\mathrm{spec}(M)\setminus\{0\}
\qquad(\text{重数相同}),
\]
则存在常数 $C\neq 0$ 使得完成化函数
\[
\Xi_M(s):=\lambda^{-rs/2}\,Z_M(s)
\]
满足函数方程
\[
\Xi_M(s)=C\,\Xi_M(1-s).
\]
\end{proposition}
\begin{proof}
按非零特征值分解
\(
\Delta_M(z)=\prod_{\mu\neq 0}(1-\mu z)
\)。
在 $s\mapsto 1-s$ 下，
\[
\Delta_M(\lambda^{-(1-s)})
=\prod_{\mu\neq 0}\bigl(1-\mu\,\lambda^{s-1}\bigr)
=\prod_{\mu\neq 0}\bigl(1-(\mu/\lambda)\,\lambda^s\bigr).
\]
用互反闭包把指标重标为 $\mu\mapsto \lambda/\mu$，并将出现的 $\lambda^{\pm rs/2}$ 幂吸收到完成化因子中，即得
\(
\lambda^{-rs/2}\Delta_M(\lambda^{-s})
=C^{-1}\lambda^{-r(1-s)/2}\Delta_M(\lambda^{-(1-s)})
\)；
取倒数即得结论。证毕。
\end{proof}

\begin{remark}[等价的 \texorpdfstring{$\Delta_M$}{Delta}-互反多项式判据]\label{rem:delta-reciprocal-polynomial}
命题 \ref{prop:reciprocal-spectrum-functional-equation} 的谱条件等价于 $z$-根集合在
\(
z\mapsto (\lambda z)^{-1}
\)
下按重数不变，亦等价于存在 $C'\neq 0$ 使得
\[
\lambda^{r}z^{r}\Delta_M\!\left(\frac{1}{\lambda z}\right)=C'\Delta_M(z).
\]
因此当 $\Delta_M$ 的系数可精确计算（例如整系数矩阵）时，上式提供一个逐项可审计的判别式。
\end{remark}

\begin{remark}[有限维与“无限零点”的缺口]\label{rem:finite-vs-infinite-zeros}
由于 \(Z_M(s)\) 来自有限维行列式，其 \((s)\)-平面的极点/零点结构必为有限族的竖向晶格（每个非零特征值贡献一条周期列）。因此若要在本框架内逼近经典 \(\zeta(s)\) 的“无限多非平凡零点”现象，需要引入一个维度 \(\to\infty\) 的极限机制，使得谱点在极限中形成无限谱并在 Dirichlet 坐标下产生更复杂的零极点分布；这为后续构造“无限塔/无限覆盖/逆极限核”的路线提供了明确的结构目标。
\end{remark}

\begin{remark}[本节两类核的检验结论（可复现）]\label{rem:reciprocal-check-outcome}
上述 $s$-平面对称并非自动成立：它要求非零谱在 $\mu\mapsto\lambda/\mu$ 下封闭。
对 10 状态同步核 $B$，由 \(\det(I-zB)\) 的因子 \((z+1)\) 知 \(\mu=-1\in\mathrm{spec}(B)\)，但 \(\lambda/\mu=-3\notin\mathrm{spec}(B)\)，故互反闭包失败。
对 40 状态真实输入核 $M$，定理 \ref{thm:pom-rh-sharp} 给出 \(\lambda^{-1}\in\mathrm{spec}(M)\setminus\{0\}\)，但 \(\lambda/(\lambda^{-1})=\lambda^2\notin\mathrm{spec}(M)\setminus\{0\}\)，故亦失败。
相应的谱互反闭包/系数互反判别由脚本 \texttt{scripts/exp\_reciprocal\_spectrum\_functional\_equation\_check.py} 自动输出到
\texttt{artifacts/export/reciprocal\_spectrum\_functional\_equation\_check.json}。
\end{remark}
